\section{Practical Diagnostics}
\label{sec:diagnostics}

The theorem's probability statement is only as meaningful as its assumptions.
Before issuing a certificate, we recommend the following audit.

\paragraph{Check support, not only average fit.}
Inspect the distribution of $\wh w(X)$ on held-out source inputs and the fraction
clipped at $B$.  Frequent clipping suggests that the numerical guarantee is tied to
a poor approximation of $w$; this should increase $\varepsilon_w$ or stop the
analysis.  Compare source and target representations using domain-relevant plots
and two-sample diagnostics \citep{gretton2012kernel}.

\paragraph{Probe conditional stability.}
Covariate shift asserts equality of $P(Y\mid X)$ and $Q(Y\mid X)$, which cannot in
general be verified from unlabeled target inputs.  Small, carefully sampled target
label audits are more informative than elaborate marginal tests.  If conditional
shift is plausible, add a separately justified sensitivity budget
$\varepsilon_{\mathrm{cond}}$ to the right side of \cref{eq:certificate}.

\paragraph{Preserve the holdout.}
The model family, loss, clipping level, and bound variant must be fixed before
examining certification outcomes.  If the sample is reused iteratively, the simple
union bound no longer describes the search.  Fresh data, nested splitting, or tools
for adaptive data analysis are then needed.

\paragraph{Choose the loss to match the harm.}
For classification error, $\ell=\mathbf{1}\{f(X)\neq Y\}$ is immediate.  Cost-
sensitive or subgroup-weighted harms can also be normalized to $[0,1]$.  A low
average loss does not imply low risk for every subgroup; certify each prespecified
group separately and include those certificates in the multiplicity budget.
